\originalchapter{8.324 作业2题解}
本章根据 MIT OCW 8.324(2010年秋)Assignment 2 题单与官方解答中文重编,补充局部性的必要条件和计算.采用本书 $(+,-,-,-)$;本套大多数结论与度规无关,规范约定仍为 $D_\mu=\partial_\mu-\ii gA_\mu$.原题1只有商书定位;原题2--6的公开题面与子问完整.

\originalsection{原题1}
\begin{exercise}
群论,原题14分:(a) Peskin--Schroeder 15.1,仅(c)、(d);(b) Peskin--Schroeder 15.2.MIT 公开题单没有给出这些商书题目的完整陈述.
\end{exercise}
\begin{solution}
保留上述精确定位.官方解答的矩阵计算可以作为学习材料,但不足以确立原商书题面及全部条件.本次未取得可再分发的完整题面,故这两个子问不重构、不计入已完成的公开原题解答.
\end{solution}

\originalsection{原题2}
\begin{exercise}
考虑 $\Lag=-\frac14F_{\mu\nu}^aF^{a\mu\nu}+\bar\Psi(\ii\slashed D-m)\Psi$.生成元满足 $[T_a,T_b]=\ii f_{abc}T_c$,$f_{abc}$ 全反对称;物质表示中的矩阵记为 $T_a^r$.原题约定写场方程为 $D_\mu F^{\mu\nu}=J^\nu$,$F^{\mu\nu}=F^{a\mu\nu}T_a^r$,$J^\nu=g(\bar\Psi\gamma^\nu T_a^r\Psi)T_a^r$.本章在 $(+,-,-,-)$、$D=\partial-\ii gA$ 的作用量约定下,物质耦合为 $+gA_\nu^aj^{a\nu}$,变分对应 $D_\mu F^{\mu\nu}=-J^\nu$;以下按这个已明确转换的符号计算.(a) 用运动方程证明流协变守恒;(b) 直接从物质场变换推导 $J^\nu$ 的规范变换,验证场方程两端相容,不能从左端反推右端.原题12分.
\end{exercise}
\begin{solution}
(a) 对场方程作用 $D_\nu$.利用 $F^{\mu\nu}$ 的反对称性,
\[
D_\nu J^\nu=-D_\nu D_\mu F^{\mu\nu}
=-\frac12[D_\nu,D_\mu]F^{\mu\nu}
=\frac{\ii g}{2}[F_{\nu\mu},F^{\mu\nu}]=0.
\]
最后一步可以在惯性坐标逐对看:对于每个固定的 $\mu,\nu$,$F_{\nu\mu}$ 与 $F^{\mu\nu}$ 只差一个实数符号,故对易子为零.也可写成 $f_{abc}F_{\mu\nu}^bF^{c\mu\nu}=0$,因为场强的收缩对 $b,c$ 对称.这是协变守恒 $D_\nu J^\nu=0$,不是一般的 $\partial_\nu J^\nu=0$.

(b) 在 $\delta\Psi=\ii\alpha^bT_b^r\Psi$、$\delta\bar\Psi=-\ii\alpha^b\bar\Psi T_b^r$ 下直接计算,
\[
\delta J^{a\nu}=\ii g\alpha^b\bar\Psi\gamma^\nu[T_a^r,T_b^r]\Psi
=-f_{abc}\alpha^bJ^{c\nu}.
\]
矩阵形式正是 $\delta J^\nu=\ii[\alpha^bT_b^r,J^\nu]$.沿连续群变换积分得 $J^{\nu\prime}=UJ^\nu U^{-1}$;也可以用 $U^{-1}T_a^rU$ 的伴随表示矩阵与其正交性直接得到有限式.另一方面,$F'=UFU^{-1}$ 以及 $D'X'=U(DX)U^{-1}$ 独立保证 $D_\mu'F^{\mu\nu\prime}=U(D_\mu F^{\mu\nu})U^{-1}$.故两端按同一伴随表示变换,检验完成.
\end{solution}

\originalsection{原题3}
\begin{exercise}
Yang--Mills 拉格朗日量可加入 $\Lag_\theta=\frac\theta{64\pi^2}\epsilon^{\mu\nu\lambda\rho}F_{\mu\nu}^aF_{\lambda\rho}^a$.(a) 讨论它在 $C,P,T$ 下的性质;(b) 将它明确写为全导数.原题12分.
\end{exercise}
\begin{solution}
(a) 令 $E_i^a=F_{0i}^a$、$B_i^a=\frac12\epsilon_{ijk}F_{jk}^a$,则该密度正比于 $\sum_a\mathbf E^a\cdot\mathbf B^a$,比例的整体符号随 $\epsilon$ 约定而变,不影响离散对称性.$P$ 下 $\mathbf E$ 为极向量、$\mathbf B$ 为轴向量,故密度为奇;$T$ 下 $\mathbf E$ 为偶、$\mathbf B$ 为奇,仍为奇.非阿贝尔 $C$ 的安全写法是矩阵场 $A\mapsto-A^T$、$F\mapsto-F^T$;因此 $\Tr(F_{\mu\nu}F_{\lambda\rho})$ 为偶,不能简单假定每个颜色分量都单独反号.结论是 $C$ 偶、$P$ 奇、$T$ 奇,因而 $CP$ 奇而 $CPT$ 偶.

(b) 将 $F^a=\dd A^a+\frac g2f_{abc}A^b\wedge A^c$ 看作二形式.展开可知 $F^a\wedge F^a$ 中四个 $A$ 的项由 Jacobi 恒等式消失;又有
\[
\dd\bigl(A^a\wedge\dd A^a\bigr)=\dd A^a\wedge\dd A^a,
\qquad
\dd\bigl(f_{abc}A^a\wedge A^b\wedge A^c\bigr)
=3f_{abc}\dd A^a\wedge A^b\wedge A^c.
\]
由此得到
\[
F^a\wedge F^a=
\dd\left(A^a\wedge\dd A^a+\frac g3f_{abc}A^a\wedge A^b\wedge A^c\right).
\]
化回分量即
\begin{align*}
\epsilon^{\mu\nu\lambda\rho}F_{\mu\nu}^aF_{\lambda\rho}^a&=\partial_\mu K^\mu,\\
K^\mu&=4\epsilon^{\mu\nu\lambda\rho}
\left(A_\nu^a\partial_\lambda A_\rho^a+
\frac g3f_{abc}A_\nu^aA_\lambda^bA_\rho^c\right),\\
\Lag_\theta&=\frac\theta{64\pi^2}\partial_\mu K^\mu.
\end{align*}
$K^\mu$ 本身一般不规范不变,所以全导数性质不表示它在所有全局拓扑扇区都没有效应.以上沿用题单的系数;本题不据此另行推断在不同场强归一化下的 $\theta$ 周期.
\end{solution}

\originalsection{原题4}
\begin{exercise}
设一个酉矩阵群规范场满足 $F_{\mu\nu}=0$,证明它局部规范等价于零.(a) 写消去 $A_\mu$ 的群元 $U(x)$ 的偏微分方程并检验可积条件;(b) 指定 $U(x_0)=1$,沿路径 $x(s)$ 写解,解释只沿一条路径为什么不够;(c) 用固定端点的路径变形证明结果与路径无关.原题24分.
\end{exercise}
\begin{solution}
这里“局部”应落实为足够小的可缩邻域;仅仅连通并不能排除非平凡的平坦 holonomy.

(a) 在 $A_\mu'=UA_\mu U^{-1}-\frac{\ii}{g}(\partial_\mu U)U^{-1}$ 中要求 $A'=0$,得
\[
\partial_\mu U=-\ii gUA_\mu.
\]
对两个方向再微分相减,直接得到
\[
(\partial_\mu\partial_\nu-\partial_\nu\partial_\mu)U
=-\ii gU\bigl(\partial_\mu A_\nu-\partial_\nu A_\mu-\ii g[A_\mu,A_\nu]\bigr)
=-\ii gUF_{\mu\nu}=0.
\]
这正是所需的局部可积条件.

(b) 与切向量收缩得 $\dot U=-\ii gU\dot x^\mu A_\mu$.这是右乘矩阵的微分方程,解为反路径排序指数
\[
U(s)=\overline{\mathcal P}\exp\!\left(-\ii g\int_0^s A_\mu(x(t))\dot x^\mu(t)\dd t\right),
\]
其中较晚矩阵在右.等价地,它是通常左排序 Wilson 线 $V=\mathcal P\exp(\ii g\int A)$ 的逆 $V^{-1}$.一条路径只检验偏微分方程在该切向量方向的投影.若端点值与路径无关,可选择末端切向量沿任意坐标方向的路径,才得到每一个 $\partial_\mu U$.

(c) 取 $\delta x(0)=\delta x(1)=0$,对路径方程变分,
\[
\delta\dot U=-\ii g\delta U\,A_\mu\dot x^\mu
-\ii gU\bigl(\partial_\nu A_\mu\,\delta x^\nu\dot x^\mu+A_\mu\delta\dot x^\mu\bigr).
\]
令 $W=\delta U\,U^{-1}$,用 $\dot U^{-1}=\ii gA_\mu\dot x^\mu U^{-1}$ 消去第一项,可得
\[
\dot W=-\ii gU\bigl(\partial_\nu A_\mu\,\delta x^\nu\dot x^\mu
+A_\mu\delta\dot x^\mu\bigr)U^{-1}.
\]
将后一项分部积分,端点项因固定端点消失,并利用
\[
\frac{\dd}{\dd s}(UA_\nu U^{-1})
=U(\partial_\mu A_\nu-\ii g[A_\mu,A_\nu])U^{-1}\dot x^\mu,
\]
得到精确的一阶路径变形公式
\[
\delta U(1)U^{-1}(1)
=\ii g\int_0^1UF_{\mu\nu}U^{-1}\dot x^\mu\delta x^\nu\dd s=0.
\]
在可缩邻域里任意两条同端点路径可以连续互相变形,沿此变形积分即知端点解相同.若区域有不可缩闭路,该论证只证明同伦类内无关,不能消去不同同伦类间的 holonomy;这也是题目必须保留“局部”的原因.
\end{solution}

\originalsection{原题5}
\begin{exercise}
(a) 证明阿贝尔规范场若有相同场强则局部规范等价.对 $SU(2)$ 取 $T_i=\sigma_i/2$,第一组 $A_x=-gyT_3/2$、$A_y=gxT_3/2$,第二组 $A_x'=T_1$、$A_y'=T_2$,其余分量均为零.(b) 证明两者场强相同;(c) 选择规范协变局域对象,证明它们仍不规范等价.原题20分.
\end{exercise}
\begin{solution}
(a) 令 $a=A'-A$.同场强意味着 $\partial_\mu a_\nu-\partial_\nu a_\mu=0$.在可缩邻域内定义 $\chi(x)=\int_{x_0}^x a_\mu\dd x^\mu$,由闭路积分为零知路径无关,故 $a_\mu=\partial_\mu\chi$.选择 $U=\e^{\ii g\chi}$,便有 $A'=A+\partial\chi$,所以局部等价.

(b) 第一组场都沿 $T_3$,对易子为零,$F_{xy}=\partial_x(gxT_3/2)-\partial_y(-gyT_3/2)=gT_3$.第二组为常数矩阵,导数为零,而
\[
F'_{xy}=-\ii g[T_1,T_2]=-\ii g(\ii T_3)=gT_3.
\]
只有 $xy,yx$ 分量非零,两组场强完全相同.

(c) 不能只比较 $A$,因为它的变换有非齐次导数项.应比较 $D_\rho F_{xy}$,它按 $X\mapsto UXU^{-1}$ 变换.第一组的 $F$ 为常数且与 $A$ 对易,因此处处 $D_\rho F_{xy}=0$.第二组则有
\[
D_x'F_{xy}'=-\ii g[T_1,gT_3]=-g^2T_2,
\qquad D_y'F_{xy}'=-\ii g[T_2,gT_3]=g^2T_1.
\]
当 $g\ne0$ 时非零;可逆共轭不可能把零矩阵变成非零矩阵.所以不存在将两连接互相映射的局部规范变换.也可比较正定颜色迹 $\Tr((D_xF_{xy})^2)$,两者分别为零与 $g^4/2$.
\end{solution}

\originalsection{原题6}
\begin{exercise}
(a) 在轴规范 $A_z^a=0$ 中导出规范场传播子;(b) 对广义 Coulomb 规范函数 $f^a=\boldsymbol\nabla\cdot\mathbf A^a$ 求鬼场拉格朗日量与鬼传播子.原题18分.
\end{exercise}
\begin{solution}
(a) 在本书度规下,取常空间向量 $n$,$n^2=-1$,规范条件 $n\cdot A=0$.动量空间二次算符可写为
\[
K^{\mu\nu}=-k^2\eta^{\mu\nu}+k^\mu k^\nu-\xi^{-1}n^\mu n^\nu.
\]
求逆时按 $\eta_{\mu\nu}$、$k_\mu k_\nu$、$k_\mu n_\nu+n_\mu k_\nu$、$n_\mu n_\nu$ 展开.将系数代入 $K^{\mu\rho}(K^{-1})_{\rho\nu}=\delta^\mu_\nu$,然后取 $\xi\to0$,得到
\[
D_{\mu\nu}^{ab}(k)=\frac{-\ii\delta^{ab}}{k^2+\ii0}
\left[\eta_{\mu\nu}-\frac{k_\mu n_\nu+n_\mu k_\nu}{k\cdot n}
+\frac{n^2k_\mu k_\nu}{(k\cdot n)^2}\right].
\]
为检验而非只引用公式,令方括号为 $P_{\mu\nu}$.直接收缩得 $n^\mu P_{\mu\nu}=0$,$k^\mu P_{\mu\nu}=k^2[-n_\nu/(k\cdot n)+n^2k_\nu/(k\cdot n)^2]$.因此 Maxwell 算符乘 $-P/k^2$ 为 $\delta^\mu_\nu-n^\mu k_\nu/(k\cdot n)$.它并非在所有 $n\cdot A=0$ 的四向量上直接等于单位算符;还要把方程投影到受约束的切空间.记该投影 $\mathcal T^\mu{}_{\nu}=\delta^\mu_\nu-n^\mu n_\nu/n^2$,便有 $\mathcal T(1-nk/(k\cdot n))A=A$($n\cdot A=0$),多出的项完全沿已删除的 $n$ 方向.若直接去掉 $A_3$ 及其方向的方程,剩下 $\widehat\mu,\widehat\nu=0,1,2$,用 $k_\parallel^2=k^2+k_3^2$ 可验证逆矩阵的秩一修正系数恰为 $1/k_3^2$,与上式非零的三乘三块一致.$k\cdot n=0$ 的额外极点来自残余规范自由度;题单未指定其分布处方,本式仅给形式传播子,不能据此宣称环积分处方已确定.

(b) 写 $\delta A_i^a=(\mathcal D_i)^{ab}\alpha^b/g$,其中 $(\mathcal D_i)^{ab}=\delta^{ab}\partial_i+gf^{acb}A_i^c$,那么
\[
\frac{\delta f^a(x)}{\delta\alpha^b(y)}
=\frac1g\partial_i(\mathcal D_i)^{ab}\delta^4(x-y).
\]
略去行列式中无场依赖的常因子 $1/g$,用 Grassmann 场指数化.选取使自由鬼核为 $-\boldsymbol\nabla^2$ 的约定,
\[
\Lag_{\rm gh}=-\bar c^a\partial_i(\mathcal D_i)^{ab}c^b
=\partial_i\bar c^a\partial_i c^a
+gf^{acb}(\partial_i\bar c^a)A_i^c c^b,
\qquad D_{\rm gh}^{ab}(k)=\frac{\ii\delta^{ab}}{\mathbf k^2}.
\]
另一种鬼场整体符号可由 $\bar c\mapsto-\bar c$ 互相转换.传播子与 $k^0$ 无关,显示 Coulomb 鬼核是瞬时的;$\mathbf k=0$ 的残余零模还需边界条件处理.
\end{solution}
